\documentclass[aspectratio=169]{beamer}
\input{header}
\coursecode{JKSSB}

\title{Types of Computers}
\subtitle{Lesson 02 --- Classify by Signal and Scale}
\author{JKSSB Computer Awareness}
\date{\today}

\begin{document}
\frame{\titlepage}

\begin{frame}{Video Lesson 02}
  \centering
  \Large\href{https://youtu.be/d6zL3YS69kM}{Watch: Computer Classifications for JKSSB}
\end{frame}

\begin{frame}{One Device, More Than One Classification}
  \begin{itemize}
    \item A laptop is a personal computer by intended use and a digital
      computer by the way it represents data.
    \item These labels answer different questions.
  \end{itemize}
  \begin{block}{Classification rule}
    First identify the basis: signal handled, size/capability, intended
    purpose, or form factor.
  \end{block}
\end{frame}

\begin{frame}{By the Kind of Data Signal}
  \begin{columns}[T]
    \column{0.31\textwidth}
      \textbf{Analog}
      \begin{itemize}
        \item Represents continuously varying quantities.
        \item Example: a traditional analog speedometer.
      \end{itemize}
    \column{0.31\textwidth}
      \textbf{Digital}
      \begin{itemize}
        \item Represents data as discrete values, commonly binary digits.
        \item Example: a laptop.
      \end{itemize}
    \column{0.31\textwidth}
      \textbf{Hybrid}
      \begin{itemize}
        \item Combines analog input or measurement with digital processing.
        \item Example: a digital patient monitor.
      \end{itemize}
  \end{columns}
\end{frame}

\begin{frame}{Analog, Digital, Hybrid: The Contrast}
  \begin{itemize}
    \item \key{Analog}: values vary continuously over a range.
    \item \key{Digital}: values are represented as distinct symbols or states.
    \item \key{Hybrid}: a system uses both forms in different stages.
  \end{itemize}
  \begin{alertblock}{Common trap}
    A digital computer can process a measurement that began as an analog
    signal; that does not make the whole computer an analog computer.
  \end{alertblock}
\end{frame}

\begin{frame}{By Traditional Size and Capability}
  \begin{itemize}
    \item \key{Microcomputer}: personal-use system, such as a desktop, laptop,
      tablet, or smartphone.
    \item \key{Minicomputer}: traditional mid-range, multi-user class; the
      label is historical and less common in modern product marketing.
    \item \key{Mainframe}: high-volume transaction and input/output workloads
      for many users.
    \item \key{Supercomputer}: very high-performance parallel computation for
      large scientific or engineering workloads.
  \end{itemize}
\end{frame}

\begin{frame}{Mainframe Is Not Supercomputer}
  \begin{columns}[T]
    \column{0.46\textwidth}
      \textbf{Mainframe}
      \begin{itemize}
        \item Many concurrent users or transactions
        \item Reliability and throughput
        \item Example workload: large banking records
      \end{itemize}
    \column{0.46\textwidth}
      \textbf{Supercomputer}
      \begin{itemize}
        \item Huge calculations that can be parallelized
        \item High computational performance
        \item Example workload: climate simulation
      \end{itemize}
  \end{columns}
  \vspace{0.3cm}
  \muted{Do not decide by physical size alone; classify by the system's
  characteristic workload.}
\end{frame}

\begin{frame}{By Purpose and Form}
  \begin{itemize}
    \item \key{General-purpose}: runs many kinds of programs, such as a
      personal computer.
    \item \key{Special-purpose}: designed for a narrower task, such as an
      embedded controller.
    \item \key{Embedded computer}: built into a larger device to control or
      support its function.
    \item Desktop, laptop, tablet, and smartphone describe common form factors,
      not analog/digital signal classes.
  \end{itemize}
\end{frame}

\begin{frame}{Quick Matching}
  \begin{itemize}
    \item Continuous temperature measurement: \key{analog signal}.
    \item Smartphone: \key{digital microcomputer}.
    \item Many bank transactions at once: \key{mainframe workload}.
    \item Parallel weather-model calculation: \key{supercomputer workload}.
    \item Washing-machine controller: \key{embedded, special-purpose}.
  \end{itemize}
\end{frame}

\begin{frame}{Generation Link: Keep the Axes Separate}
  \begin{itemize}
    \item Computer \key{type} can mean its signal class, scale, purpose, or
      form.
    \item Computer \key{generation} refers to historical technology eras.
    \item The next history lesson connects vacuum tubes, transistors,
      integrated circuits, and microprocessors to generations.
  \end{itemize}
  \begin{block}{Exam cue}
    Integrated circuits are associated with the third generation; the
    microprocessor is associated with the fourth.
  \end{block}
\end{frame}

\begin{frame}{Lesson 02: Recall}
  \begin{itemize}
    \item Analog = continuous; digital = discrete; hybrid = both.
    \item Microcomputer = personal-use category; embedded = built into a
      larger product.
    \item Mainframe emphasizes throughput and many users; supercomputer
      emphasizes high-performance computation.
    \item Name the basis of classification before selecting an answer.
  \end{itemize}
\end{frame}
\end{document}
